% !TeX program = xelatex
% disertasi.tex — Template XeLaTeX (versi bersih & rapi)
\documentclass[12pt,a4paper,oneside]{report}

% ---------------- Layout & Paket Dasar ----------------
\usepackage[a4paper, top=4cm, left=3cm, right=3cm, bottom=3cm]{geometry}
% \usepackage{showframe}    % untuk menampilkan margin (bisa dihapus nanti)
\usepackage{setspace}      % untuk kontrol spasi
\usepackage{titlesec}      % untuk atur ukuran heading
\usepackage{float}         % untuk kontrol posisi gambar/tabel
% Reduce top space before chapter title
\titlespacing*{\chapter}{0pt}{-20pt}{20pt}
\usepackage{graphicx}      % untuk gambar
\usepackage{listings}      % untuk "Daftar Tampilan" (kode/program)
\usepackage{caption}       % kontrol caption
\usepackage{indentfirst}   % ✅ buat paragraf pertama setiap bab ikut menjorok
\usepackage{ragged2e}      % ✅ menyediakan \justifying

% ---------------- Bahasa & Font (XeLaTeX) --------------
\usepackage{polyglossia}
\setmainlanguage{bahasai}  % ✅ gunakan bahasa Indonesia (tanpa warning)
\usepackage{fontspec}
\setmainfont{Times New Roman} % ganti ke TeX Gyre Termes bila TNR tidak tersedia

% ---------------- Penamaan elemen Indonesia ------------
\addto\captionsbahasai{%
  \renewcommand{\contentsname}{Daftar Isi}%
  \renewcommand{\listfigurename}{Daftar Gambar}%
  \renewcommand{\listtablename}{Daftar Tabel}%
  \renewcommand{\lstlistlistingname}{Daftar Tampilan}%
  \renewcommand{\abstractname}{Abstrak}%
}

% ---------------- Ukuran heading -----------------------
% Judul Bab = 14pt bold
\titleformat{\chapter}
  {\bfseries\fontsize{14pt}{16pt}\selectfont}
  {\thechapter}{1em}{}
% Subjudul (section, subsection) = 13pt bold
\titleformat{\section}
  {\bfseries\fontsize{13pt}{15pt}\selectfont}
  {\thesection}{1em}{}
\titleformat{\subsection}
  {\bfseries\fontsize{13pt}{15pt}\selectfont}
  {\thesubsection}{1em}{}

% ---------------- Bibliografi (spasi 1) ----------------
\usepackage[backend=biber,style=ieee,sorting=nyt,autolang=other]{biblatex}
\DeclareLanguageMapping{bahasa}{english}
\DeclareLanguageMapping{bahasai}{english}
\addbibresource{references.bib}
\AtBeginBibliography{\singlespacing}

% ---------------- Listings (opsional) -------------------
\lstset{
  basicstyle=\ttfamily\small,
  numbers=left,
  numberstyle=\tiny,
  stepnumber=1,
  numbersep=5pt,
  frame=single,
  breaklines=true
}

% =======================================================
\begin{document}

  % ===================== HALAMAN JUDUL (spasi 1) =====================
  \begin{titlepage}
      \thispagestyle{empty}
      \begin{center}
          \Large{\textbf{KECERDASAN BUATAN LANJUT}}\\[0.5cm]
          \Large{\textbf{TUGAS KULIAH}}\\[0.5cm]

          \large{\textbf{disusun oleh}}\\
          \large{\textbf{Benny Leonard Enrico Panggabean}}\\
          \large{\textbf{D063251008}}\\[2cm]

          \includegraphics[width=5cm]{Logo-Resmi-Unhas-1.png}\\[1.5cm]

          \large{\textbf{PROGRAM STUDI DOKTOR TEKNIK INFORMATIKA}}\\
          \large{\textbf{FAKULTAS TEKNIK}}\\
          \large{\textbf{UNIVERSITAS HASANUDDIN}}\\
          \large{\textbf{MAKASSAR}}\\
          \large{\textbf{2025}}\\
      \end{center}
  \end{titlepage}
  % ===================== TEKS UTAMA (spasi 1.5) ======================
  \onehalfspacing
  \justifying
  \chapter{Pendahuluan}
    Tugas Kuliah disusun sesuai dengan arahan yang diberikan oleh dosen pemateri perkuliahan,
    dimana tugas ini menjadikan topik tugas akhir sebagai bahan kajian dan pembahasan dalam penyusunan tugas kuliah ini.
    Draft tugas akhir yang disusun pada tugas kuliah ini masih bersifat sementara dan akan terus dikembangkan
    sehingga menjadi tugas akhir yang lengkap dan utuh. Judul tugas akhir yang diambil adalah
    \textit{“DATA ALIGNMENT BERBASIS PROBABILISTIC ENTITY RESOLUTION DAN DETEKSI ANOMALI KEPATUHAN YANG TERINTEGRASI DENGAN RETRIEVAL AUGMENTED GENERATION.”}
    Secara ringkas, alur penelitian dapat dijelaskan melalui konsep berikut:
    \vspace{0.5\baselineskip}
    Data SIAKAD dan FEEDER $\rightarrow$ Entity Resolution $\rightarrow$ Anomaly Detection $\rightarrow$ RAG (QA berbukti regulasi + data).
    \vspace{0.5\baselineskip}
    secara umum gambar kerangka konseptual dapat dilihat pada Gambar diagaram pada halaman berikutnya,
    dapat dijelaskan bahwa data SIAKAD dan FEEDER yang berasal dari sumber berbeda
    akan dilakukan proses \textit{Entity Resolution} untuk menyelaraskan data tersebut.
    Setelah data selaras, dilakukan proses \textit{Anomaly Detection} untuk mendeteksi ketidakpatuhan data
    terhadap regulasi yang berlaku. Akhirnya, data yang telah melalui proses-proses tersebut
    akan digunakan dalam sistem \textit{Retrieval Augmented Generation} (RAG)
    untuk menjawab pertanyaan-pertanyaan terkait regulasi dan data yang ada.
    \par
      tahap pertama yang dilakukan adalah pengambilan data dari FEEDER DIKTI melalui API yang disediakan.
      Data yang diambil meliputi data mahasiswa, data dosen, data mata kuliah, data nilai, dan data lainnya yang relevan.
      Data ini kemudian akan disimpan dalam basis data lokal untuk proses selanjutnya.
    \par 
      \begin{figure}[H]
          \centering
          \includegraphics[width=0.7\linewidth]{DiagramKerangkaKonseptual.png}
          \caption*{Diagram Kerangka Konseptual}
      \end{figure}

      \vspace{0.5\baselineskip}
      \justifying

      Proses pengambilan data dari FEEDER DIKTI melalui API dilakukan dengan menggunakan
      bahasa pemrograman Rust, dengan menggunakan metode paginasi untuk mengambil data
      dalam jumlah besar. Data yang diambil kemudian disimpan dalam basis data lokal
      menggunakan RDBMS PostgreSQL. Proses ini memastikan bahwa data yang diambil konsisten dengan data yang ada di FEEDER DIKTI dan yang diambil.
    \par
      proses pengambilan data terjadi dua proses, proses pertama adalah proses pengambilan data dengan pengambilan data sebanyak 1000 data per halaman, data tersebut disimpan
      dalam proses antrian (queue) untuk diproses pada tahap selanjutnya.
      Proses kedua adalah proses penyimpanan data dari antrian ke basis data lokal PostgreSQL, proses ini melakukan proses upsert data, dimana jika data sudah ada maka data akan diperbarui,
      jika data belum ada maka data akan disisipkan sebagai data baru.
      Proses ini memastikan bahwa data yang disimpan dalam basis data lokal selalu up-to-date dengan data yang ada di FEEDER DIKTI. 

      \begin{figure}[H]
          \centering
          \includegraphics[width=0.7\linewidth]{task.png}
          \caption*{Task}
      \end{figure}

    \par
      gambar diatas menunjukkan proses pengambilan data dari FEEDER DIKTI melalui API,
      dimana proses pengambilan data dilakukan secara bertahap dengan menggunakan metode paginasi.
      Data yang diambil kemudian disimpan dalam basis data lokal PostgreSQL untuk proses selanjutnya.
      Proses ini memastikan bahwa data yang diambil konsisten dengan data yang ada di FEEDER DIKTI dan yang diambil.
      sedangkan gambar berikut adalah gambar ketika proses upsert data ke basis data lokal PostgreSQL.\
      \begin{figure}[H]
          \centering
          \includegraphics[width=0.7\linewidth]{worker.png}
          \caption*{Worker}
      \end{figure}
  \chapter{Penutup}
    Tugas kuliah ini merupakan bagian dari proses penyusunan tugas akhir yang berjudul
    \textit{“DATA ALIGNMENT BERBASIS PROBABILISTIC ENTITY RESOLUTION DAN DETEKSI ANOMALI KEPATUHAN YANG TERINTEGRASI DENGAN RETRIEVAL AUGMENTED GENERATION.”}
    dimana tugas kuliah ini membahas tentang pengambilan data dari FEEDER DIKTI melalui API
    dan penyimpanan data dalam basis data lokal PostgreSQL.
    Proses pengambilan data dilakukan dengan menggunakan bahasa pemrograman Rust
    dan metode paginasi untuk mengambil data dalam jumlah besar.
    Data yang diambil kemudian disimpan dalam basis data lokal menggunakan RDBMS PostgreSQL.
    Proses ini memastikan bahwa data yang diambil konsisten dengan data yang ada di FEEDER DIKTI dan yang diambil.
    Tugas kuliah ini akan terus dikembangkan sebagai bagian dari penyusunan tugas akhir yang lengkap dan utuh.
\end{document}
